El presente capítulo establece las políticas, procedimientos y responsabilidades para la gestión, control, actualización y distribución de la información técnica utilizada en las actividades de mantenimiento de los Sistemas de Aeronaves No Tripuladas (UAS) de SKYFILMS.

La Biblioteca Técnica tiene como propósito garantizar que todo el personal técnico disponga de documentación vigente, confiable y verificada, asegurando que el mantenimiento se realice conforme a los Reglamentos Aeronáuticos de Colombia (RAC 100), los manuales del fabricante, las directrices de la Aerocivil y las mejores prácticas aeronáuticas aplicables a operaciones con UAS.

\section{Alcance}
Este procedimiento aplica a todos los documentos técnicos que se utilicen en la ejecución del mantenimiento preventivo, correctivo y de inspección de los UAS de SKYFILMS, incluyendo pero sin limitarse a:

\begin{itemize}
    \item Manuales de operación y mantenimiento de cada modelo de UAS.

    \item Boletines de servicio, avisos técnicos y actualizaciones del fabricante.

    \item Programas y procedimientos de mantenimiento aprobados por SKYFILMS.

    \item Fichas técnicas de componentes, sensores y baterías.

    \item Normas, circulares y documentos técnicos emitidos por la Aerocivil.

    \item Procedimientos internos y formatos aprobados del MCM.
\end{itemize}

\section{Organización de la Biblioteca Técnica}
La Biblioteca Técnica de SKYFILMS está organizada y administrada por el \textbf{Técnico de UAS},y cuenta con los siguientes componentes:

\begin{enumerate}
    \item \textbf{Archivo físico:} ubicado en la \textbf{Oficina Técnica de Mantenimiento}, con copias impresas controladas de los documentos esenciales, disponibles para consulta inmediata.

    \item \textbf{Archivo digital:} repositorio electrónico centralizado en el sistema técnico de SKYFILMS, con control de acceso restringido, copia de respaldo automática y trazabilidad de versiones.

    \item \textbf{Base de datos de control documental:} registro actualizado de cada documento técnico, indicando número de versión, fecha de emisión, origen y estado de vigencia.
\end{enumerate}

Todo documento contenido en la Biblioteca Técnica debe estar identificado, numerado y controlado según el procedimiento de gestión documental establecido por SKYFILMS.

\section{Documentos del Fabricante por Modelo}
Para los modelos de la flota, los fabricantes (DJI y BetaFPV) publican manuales de usuario, guías de seguridad de baterías, notas de versión de firmware y avisos de seguridad. Hasta donde Skyfilms ha verificado en sus sitios oficiales, no publican catálogo ilustrado de partes (IPC) ni manual de taller para estos equipos de uso general, por lo que el mantenimiento se rige por los manuales de usuario. Si alguno los publica después, el Técnico UAS los incorpora a la biblioteca.

\begin{longtable}{|>{\raggedright\arraybackslash}p{4.6cm}|>{\raggedright\arraybackslash}p{3.2cm}|>{\raggedright\arraybackslash}p{5.9cm}|}
\hline
\rowcolor[gray]{0.9}\textbf{Documento} & \textbf{Modelo} & \textbf{Uso en mantenimiento} \\ \hline
\endhead
Manual de usuario DJI Mavic 4 Pro & Mavic 4 Pro & Inspecciones, limpieza, hélices y procedimientos de uso \\ \hline
Guía de seguridad de baterías de vuelo inteligentes DJI & Mavic 4 Pro & Límites y cuidado de baterías \\ \hline
Manual de usuario DJI RC Pro 2 & Mavic 4 Pro & Control remoto \\ \hline
Manual y guía de BetaFPV Pavo20 Pro II & Pavo20 Pro II & Estructura, motores, hélices y ajustes \\ \hline
Documentación de DJI O4 Pro, ExpressLRS, Betaflight y GPS & Pavo20 Pro II & Enlace de video y control, configuración y navegación \\ \hline
Notas de versión de firmware & Ambos & Control de versiones en el Libro de Mantenimiento \\ \hline
Boletines o avisos de seguridad del fabricante & Ambos & Revisión semestral (ver abajo) \\ \hline
\end{longtable}

Cada vez que el Técnico UAS revisa las fuentes oficiales, anota en el inventario de la biblioteca (base de datos de control documental) la fecha de la consulta, el documento, su versión y si hay cambios. Esa tabla es el registro de consulta de publicaciones técnicas y el inventario de la biblioteca.

\section{Control y Actualización de Documentos Técnicos}

\begin{itemize}
    \item Solo se consideran válidos los documentos emitidos o aprobados por el fabricante del UAS o por la Aerocivil.

    \item El \textbf{Técnico UAS} es responsable de revisar \textbf{semestralmente} las fuentes oficiales del fabricante y las circulares de la autoridad para detectar nuevas versiones, boletines o revisiones técnicas.

    \item Cualquier documento nuevo o actualizado debe ser evaluado y aprobado por el \textbf{Técnico UAS} antes de su incorporación al sistema.

    \item Una vez aprobada, la nueva versión reemplaza inmediatamente a la anterior, la cual debe ser retirada del uso operativo y archivada en la sección de documentos históricos.

    \item Cada documento actualizado debe reflejar su \textbf{número de revisión y fecha de emisión} en la portada o pie de página, conforme al formato de control interno.
\end{itemize}

La Biblioteca Técnica mantiene además un registro de control de enmiendas donde se documenta cada actualización, su fuente, responsable y fecha de implementación.

\section{Distribución y Acceso}
\begin{itemize}
    \item El personal técnico de mantenimiento debe tener acceso garantizado a la documentación técnica vigente correspondiente al modelo de UAS que opere o mantenga.

    \item El acceso se realiza mediante el sistema digital autorizado o las copias controladas disponibles en la \textbf{Oficina Técnica}.

    \item \textbf{Queda prohibida la reproducción o distribución de documentos técnicos sin autorización del Técnico UAS}, a fin de evitar el uso de versiones no controladas.

    \item En caso de mantenimiento fuera de la base principal, el personal técnico deberá portar copias electrónicas o impresas actualizadas y verificadas del material técnico necesario para la tarea específica.
\end{itemize}

\section{Control de Documentos Obsoletos}
\begin{itemize}
    \item Todo documento técnico que sea reemplazado o declarado obsoleto deberá ser retirado inmediatamente de las áreas operativas y marcado con el sello \textbf{“NO VIGENTE”} o su equivalente digital.

    \item Los documentos obsoletos se conservarán únicamente con fines de trazabilidad histórica y permanecerán archivados por un \textbf{periodo mínimo de cinco (5) años} o según lo determine la Aerocivil.

    \item El \textbf{Técnico UAS} es responsable de garantizar que no exista material técnico obsoleto en uso dentro de las áreas de mantenimiento o vuelo.
\end{itemize}

\section{Responsabilidades}

\begin{itemize}
    \item El \textbf{Técnico UAS} es:
        \begin{itemize}
            \item Responsable de la supervisión general del control y actualización de la Biblioteca Técnica.
            
            \item Responsable de administrar los registros, gestiona la actualización documental y asegura la disponibilidad y accesibilidad del material técnico.

           \item Responsable de consultar exclusivamente la documentación vigente antes de iniciar cualquier tarea, y reportar si se identifican discrepancias o desactualizaciones.

            \item Responsable de garantizar los recursos tecnológicos y logísticos para mantener la Biblioteca Técnica actualizada y disponible en todo momento.
        \end{itemize}
    
\end{itemize}

\section{Disposiciones Finales}
El mantenimiento de la Biblioteca Técnica constituye un requisito esencial para asegurar la validez y conformidad del programa de mantenimiento ante la Unidad Administrativa Especial de Aeronáutica Civil (Aerocivil).\\

Toda auditoría, inspección o revisión técnica podrá requerir la evidencia de los documentos vigentes y su trazabilidad dentro del sistema.\\

El incumplimiento de los procedimientos aquí establecidos será considerado una no conformidad técnica mayor, sujeta a acciones correctivas inmediatas.\\

El control riguroso de la Biblioteca Técnica garantiza que SKYFILMS ejecute todas las actividades de mantenimiento de sus UAS bajo los más altos estándares de precisión técnica, seguridad operacional y cumplimiento normativo.